% l1_content.tex
% Частина конспекту лекції "Python. Вступ та синтаксис", оформлена відповідно до вимог ОНУ.





\chapter{\parbox{0.85\textwidth}{Чому Python?  Створення та запуск файлів. Робота в інтерпретаторі.}}

\section{Історія та Глобальний вплив}

Створення Python було розпочато Гвідо ван Россумом ( Guido van Rossum ) у 1991 році, коли він працював над розподіленою ОС Амеба. Йому була потрібна мова, що розширюється, який би забезпечив підтримку системних викликів. За основу були взяті ABC та Модула-3. Як назви він вибрав Python на честь комедійних серій BBC "Літаючий цирк Монті-Пітона", а зовсім не за назвою змії. З того часу Python розвивався за підтримки тих організацій, у яких Гвідо працював. Особливо активно мову вдосконалюється нині, коли з неї працює як команда творців, а й ціле співтовариство програмістів з усього світу. І все-таки останнє слово про напрям розвитку мови залишається за Гвідо ван Россумом.

\subsection{Історія розвитку Python}

\begin{itemize}
	\item \textbf{Python 2.0} --- випущений у 2000 році. Серед нововведень: спискові включення, підтримка Unicode, вдосконалений збирач сміття.
	\item \textbf{Python 3.0} --- випущений у 2008 році. Основна особливість --- несумісність із попередніми версіями, що вимагало адаптації бібліотек.
	\item \textbf{Python 2.7} --- останній реліз гілки 2.x. Його підтримку офіційно завершено у 2020 році, що ознаменувало перехід до епохи Python~3.
\end{itemize}

\subsection{Глобальний вплив та переваги}

Python став однією з найпопулярніших мов програмування у світі завдяки простоті синтаксису та універсальності. Його активно використовують такі компанії, як \textbf{Google, Netflix, Facebook, Instagram, Spotify} та багато інших.

\medskip
Основні сфери застосування:
\begin{itemize}
	\item \textbf{Наука та дослідження} --- бібліотеки \texttt{NumPy}, \texttt{SciPy}, \texttt{pandas}.
	\item \textbf{Машинне навчання та штучний інтелект} --- \texttt{TensorFlow}, \texttt{PyTorch}, \texttt{scikit-learn}.
	\item \textbf{Веб-розробка} --- фреймворки \texttt{Django}, \texttt{Flask}.
	\item \textbf{Автоматизація та DevOps} --- скрипти для адміністрування, тестування, CI/CD.
	\item \textbf{Аналіз даних та візуалізація} --- \texttt{matplotlib}, \texttt{seaborn}, \texttt{Jupyter Notebooks}.
\end{itemize}

\medskip
Таким чином, Python використовується практично всюди --- від навчальних проєктів до масштабних систем світових технологічних гігантів.



\section{Переваги мови або чому Python}

\subsection{Основні переваги}

\begin{itemize}
	\item \textbf{Якість програмного забезпечення.} Зручність читання, ясність і висока якість відрізняють Python від інших мов. Єдиноманітність оформлення коду полегшує його розуміння навіть для тих, хто не брав участі у створенні програми.
	
	\item \textbf{Висока швидкість розробки.} У порівнянні з C, C++ та Java, Python значно підвищує продуктивність праці розробника. Обсяг коду зазвичай становить третину або навіть п’яту частину еквівалентного коду на C++ чи Java. Програми запускаються відразу, минаючи етапи компіляції та зв’язування.
	
	\item \textbf{Перенесення програмного коду без витрат.} Більшість програм виконується без змін на основних платформах. Перенесення з Linux у Windows часто зводиться до копіювання файлів. Python також підтримує створення графічних інтерфейсів, доступ до баз даних, веб-додатків та інших типів програм.
	
	\item \textbf{Бібліотеки підтримки.} У складі Python є велика стандартна бібліотека, що охоплює пошук тексту, роботу з мережею та інші функції. Крім того, існують численні сторонні бібліотеки: для веб-розробки, математичних обчислень, ігор тощо. Наприклад, \texttt{NumPy} --- потужний інструмент для наукових обчислень.
	
	\item \textbf{Інтеграція компонентів.} Python легко взаємодіє з іншими мовами та платформами: може викликати функції з бібліотек на C/C++, інтегруватися з Java, взаємодіяти з COM та .NET, працювати через протоколи SOAP, XML-RPC, CORBA.
	
	\item \textbf{Вільне використання.} Python можна використовувати та розповсюджувати абсолютно безкоштовно.
\end{itemize}

\subsection{Класифікація мови}

\begin{enumerate}
	\item \textbf{Високого рівня.} Код легко читати і писати, він схожий на англійську мову.
	\item \textbf{Загального призначення.} Python вирішує широкий спектр завдань.
	\item \textbf{Інтерпретована.} Для виконання коду не потрібна компіляція.
\end{enumerate}

Python часто позиціонують як \textbf{об’єктно-орієнтовану мову сценаріїв}, що дозволяє створювати прості програми у вигляді одного файлу, використовуючи об’єкти.

\subsection{Недоліки}

\begin{itemize}
	\item \textbf{Швидкодія.} Програми на Python можуть працювати повільніше за аналоги, написані компілюючими мовами (наприклад, C).
	\item \textbf{Байт-код.} Python перетворює вихідний код у байт-код, який інтерпретується. Це забезпечує переносимість програм, але відсутність машинного коду може знижувати продуктивність.
\end{itemize}

\section{Середовище для роботи та Інтерпретатор}


Python --- це не лише мова програмування, а й програмний пакет, що включає \textbf{інтерпретатор} --- модуль, який створює та виконує інші програми.

\subsection{Онлайн середовища розробки}

Сьогодні є можливість писати код на Python у веб-сервісах, які дозволяють запускати програми прямо у браузері, без встановлення додаткового програмного забезпечення:
\begin{itemize}
	\item \url{https://replit.com/~}
	\item \url{https://ideone.com/}
	\item \url{https://www.online-python.com/}
\end{itemize}

\medskip
У цих сервісах достатньо створити акаунт, написати код та запустити його онлайн.

\subsection{Командний рядок інтерпретатора}

Найпростіший спосіб писати й запускати програми на Python --- вводити інструкції безпосередньо у \textbf{командний рядок інтерпретатора}. Це дозволяє швидко перевіряти окремі команди та експериментувати з кодом.

\subsection{Локальне встановлення Python}

Щоб працювати локально, потрібно завантажити Python з офіційного сайту:
\begin{itemize}
	\item \url{https://www.python.org/downloads/}
\end{itemize}

Після встановлення пакету Python створюється інтерпретатор та бібліотека підтримки. Також доступне інтегроване середовище розробки \textbf{IDLE}, яке встановлюється разом із Python і дозволяє писати, запускати та відлагоджувати програми у зручному графічному інтерфейсі.

\section{Виконання вказівок у Python}

\subsection{Інтерактивний режим}

В інтерактивному режимі користувач вводить інструкції та натискає \texttt{Enter}, після чого результат відображається рядком нижче.

\begin{verbatim}
	>>> 2 + 2
	4
	>>> "Hello"
	'Hello'
\end{verbatim}

Інтерактивний режим дозволяє:
\begin{itemize}
	\item швидко перевіряти фрагменти коду;
	\item проводити експерименти без створення файлів;
	\item одразу виявляти помилки.
\end{itemize}

Приклад помилки:
\begin{verbatim}
	>>> print(y)
	Traceback (most recent call last):
	File "<pyshell#15>", line 1, in <module>
	print(y)
	NameError: name 'y' is not defined
\end{verbatim}

\subsection{Виведення повідомлень на екран}

Для виведення інформації використовується команда \texttt{print}:

\begin{verbatim}
	>>> print("Hello World")
	Hello World
	>>> print(2 + 2)
	4
\end{verbatim}

У інтерактивному режимі можна вводити вирази без \texttt{print}, але у файлах інструкцію \texttt{print} потрібно писати обов’язково.

\subsection{Програми у файлах}

Програму на Python можна написати у будь-якому текстовому редакторі та зберегти у файл з розширенням \texttt{.py}.  
Наприклад, створимо файл \texttt{hello.py}:

\begin{verbatim}
	print("Hello world!")
	print("Привіт світ")
	input()
\end{verbatim}

Запуск програми здійснюється подвійним клацанням по файлу.  
Такі програми виконуються у консольному вікні («чорному вікні»), де введення здійснюється з клавіатури. Якщо у програмі є помилки, у консолі буде показано повідомлення про них.

\medskip
Таким чином, Python дозволяє працювати як у інтерактивному режимі, так і створювати повноцінні програми у файлах.

\section{Перевірка коду на помилки в IDLE}

\subsection{Запуск програми у IDLE}

\begin{enumerate}
	\item Запускаємо середовище \textbf{IDLE}.
	\item У меню \textbf{File} обираємо пункт \textbf{Open} та у вікні шукаємо файл \texttt{hello.py}.
	\item Відкриється нове вікно з текстом програми. У цьому вікні у меню \textbf{Run} вибираємо \textbf{Run Module} або натискаємо клавішу \texttt{F5}.
\end{enumerate}

\subsection{Результати виконання}

\begin{itemize}
	\item Якщо у програмі немає помилок, введення та виведення здійснюється у вікні інтерпретатора.
	\item Повідомлення про помилки також відображаються у \textbf{IDLE}.
	\item Якщо програма працює некоректно, але не видає помилок, для пошуку проблем можна вставити додаткові інструкції \texttt{print}.
\end{itemize}

\subsection{Приклад використання \texttt{print} для відлагодження}

\begin{verbatim}
	print("Значення змінної x:", x)
\end{verbatim}

Це допоможе зрозуміти, які значення мають змінні під час виконання програми.

\chapter{\parbox{0.84\textwidth}{Синтаксис Python. Типи Даних}}

\section{Змінні та Ідентифікатори}

\begin{defn}{Змінна - це}
	ім'я (ідентифікатор), яке може набувати певного значення.
\end{defn}

У Python змінні не потрібно оголошувати, значення можна надати одразу за допомогою символу \texttt{=}.

\begin{py}{Приклади присвоєння}
x = 4
s1 = "Hello"
_x = True
x2 = 23.46
summa1 = summa2 = k = 0
\end{py}

\begin{defn}{Ідентифікатор - це}
	унікальна ознака об'єкта, що дозволяє розрізняти його з іншими.
\end{defn}

\subsection{Правила побудови ідентифікаторів}

\begin{itemize}
	\item Використовуються лише латинські літери, цифри та знак підкреслення (\_).
	\item Ідентифікатор повинен починатися з літери або символу підкреслення.
	\item Регістр літер має значення (\texttt{X} $\neq$ \texttt{x}).
	\item Неприпустимі пробіли, точки, спеціальні символи та літери кирилиці.
	\item Не можна застосовувати спеціальні символи \texttt{\$}, \texttt{\%}, \texttt{@}.
	\item Заборонено використовувати зарезервовані слова Python (\texttt{if}, \texttt{for}, \texttt{while}, \texttt{class}, \texttt{def} тощо).
\end{itemize}

\section{Зарезервовані слова Python}

У Python існує набір \textbf{ключових слів}, які не можна використовувати як ідентифікатори (назви змінних, функцій, класів тощо).  
Ці слова мають спеціальне значення у синтаксисі мови.

\begin{center}
	\begin{tabular}{lllll}
		\texttt{False} & \texttt{class} & \texttt{finally} & \texttt{is} & \texttt{return} \\
		\texttt{None} & \texttt{continue} & \texttt{for} & \texttt{lambda} & \texttt{try} \\
		\texttt{True} & \texttt{def} & \texttt{from} & \texttt{nonlocal} & \texttt{while} \\
		\texttt{and} & \texttt{del} & \texttt{global} & \texttt{not} & \texttt{with} \\
		\texttt{as} & \texttt{elif} & \texttt{if} & \texttt{or} & \texttt{yield} \\
		\texttt{assert} & \texttt{else} & \texttt{import} & \texttt{pass} & \texttt{break} \\
	\end{tabular}
\end{center}

\medskip
\noindent Використання будь-якого з цих слів як назви змінної або функції призведе до помилки.

\subsection{Приклад помилки}

\begin{verbatim}
	>>> class = 5
	File "<stdin>", line 1
	class = 5
	^
	SyntaxError: invalid syntax
\end{verbatim}

\section{Оператори та Вирази}

\begin{defn}{Вираз — це}
	фрагмент мови програмування, що представляє спосіб обчислення деякого значення.
\end{defn}
\begin{defn}{Операнд — це}
	частина виразу, над яким виконуються операції.
\end{defn}
\begin{defn}{Оператор — це}
	дія, яка може бути виконана над одним або декількома операндами.
\end{defn}

\begin{py}{Приклади синтаксису}
X=4+5 % Присвоєння
Print(x)
S="Hello "+"World"
Print(s)
S="Helo"+\
"World" % Перенесення довгих операторів
\end{py}

\subsection{Правила форматування коду}
\begin{itemize}
	\item Кожен оператор завершується символом нового рядка.
	\item Відступи застосовуються для позначення різних блоків коду (тіла функцій, умовні висловлювання, цикли).
	\item Пробіли неприпустимі перед операторами, але необхідні перед вкладеними блоками (\texttt{unexpected indent}).
	\item Якщо потрібно розмістити кілька операторів в одному рядку, їх розділяють крапкою з комою (\texttt{;}).
	\item Символ \texttt{\#} означає коментар.
	\item Багаторядковий коментар укладають у три подвійні лапки в \texttt{py} файлі.
\end{itemize}

\section{Основні Типи Даних і Операції}

Вбудовані типи даних:
\begin{itemize}
	\item Числа: \texttt{1234}, \texttt{3.14}, \texttt{3+4j}.
	\item Рядки: \texttt{'spam'}, \texttt{"guido's"}.
	\item Списки: \texttt{[1, [2, 'three'], 4]}.
	\item Словники: \texttt{\{'food': 'spam', 'taste': 'yum'\}}.
	\item Кортежі: \texttt{(1, 'span', 4, 'U')}.
	\item Файли: myfile = open('eggs', ‘r')
	\item  Множини: set('abc'). {'a','b', 'с'}
	\end{itemize}
	
	
	\hspace{5cm}\textbf{Вбудовані числові типи}
	\begin{itemize}
	\item 
	цілі числа (числа без дробової частини), (3, 3465, 45, 42, -4) (з необмеженою точністю)
	\item дійсні числа (грубо кажучи, числа з десятковою точкою з плаваючою точкою) (3,15 5,68)
	\item  8-річні, починаються нуль про 0о77 = 63 0о10 = 8
	\item 16-річні, починаються нуль x 0xFF=255 0xff=255 • - двійкові, починаються нуль b 0b1111 = 15 0b10 = 2
	\item - комплексні числа з уявною частиною, (3+4j), (3+4J)(з допомогою великої та маленької) літер)
	\end{itemize}
	Символ множення перед j не вказується.
	\begin{pydef} 
>>>z1=2+3j
>>>z2=4+8J
>>>print(z1, z2)
(2+3j) (4+8j)
	\end{pydef}
	
	Комплексні числа можуть бути створені за допомогою функції complex(реч \_ частина, уявна \_ частина)\\
	
	\begin{pydef} 
>>>z3=complex(34,2)
>>>print(z3)
(34+2j)
	\end{pydef}
	
	\textbf{Операції з числовими типами}\\
	
	\begin{tabular}{|c|c|c|c|} 
	\hline
	\textbf{Оператор} & \textbf{Опис}&\textbf{Приклад}&\textbf{Результат}\\ \hline
	+ & Додавання&5+7&12\\ \hline
	-&Віднімання&55-10&45\\ \hline 
	*& Множення& 4*6 &24\\ \hline
	/& Ділення &7 / 2 &3.5\\ \hline
	//& Цілочисельне ділення& 7 // 2& 3\\ \hline
	\%& Остача від ділення& 7 \%  3 &1\\ \hline
	**& Піднесення до степеня& 3 ** 4 &81\\ \hline
	
	\end{tabular}
	\\
	
	Результатом операції буде новий об'єкт. \\
	 2 \bullet \bullet  100 \# 2 у степені 100\\ 
	1267650600228229401496703205376 
	
	Зверніть увагу на результат останньої операції: у Python 3 цілі числа автоматично забезпечують \textbf{необмежену} точність. Наприклад, ви можете спробувати обчислити 2 ступенем 1 000 000 (але навряд чи варто це робити, оскільки на екран буде виведено число довжиною понад 300 000 знаків, що може зайняти тривалий час!).\\
	
	\textbf{ > len(str(2 ** 1000000))} 
	
	Скільки цифр у дійсно \textsc{великому} числі\\
	301030
		
	\section{Приклади. Виділення цифр числа}
	
	\subsection{Двоцифрове число} Нехай маємо число \texttt{n = 47}.  	Щоб виділити окремі цифри:
	
	\begin{py}{Двоцифрове число}
n = 47
d1 = n // 10   # перша цифра (десятки)
d2 = n % 10    # друга цифра (одиниці)
print("Десятки =", d1, "Одиниці =", d2)
	\end{py}
	
	Результат:
	\begin{verbatim}
Десятки = 4 Одиниці = 7
	\end{verbatim}
	
	\subsection{Трицифрове число}
	
	Нехай маємо число \texttt{m = 583}.  
	Щоб виділити окремі цифри:
	
	\begin{py}{Трицифрове число}
m = 583
sot = m // 100          # сотні
des = (m // 10) % 10    # десятки
od = m % 10             # одиниці
print("Сотні =", sot, "Десятки =", des, "Одиниці =", od)
	\end{py}
	
	Результат:
	\begin{verbatim}
Сотні = 5 Десятки = 8 Одиниці = 3
	\end{verbatim}
	
	
	
	
	\emph{\textbf{Пріоритет операторів - порядок виконання  }}\href{https://pythonguide.rozh2sch.org.ua/#operators-priority}{тут}
	
	\section{Вбудовані функції для роботи з числами}
	
	Python надає низку \textbf{вбудованих функцій}, які можна застосовувати до числових типів.
	
	\subsection{Функція \texttt{abs(a)}}
	Повертає модуль числа.
	\begin{py}{Приклад}
x = -12
print(abs(x))   # 12
	\end{py}
	
	\subsection{Функція \texttt{divmod(a, b)}}
	Виконує ділення \texttt{a} на \texttt{b} та повертає пару чисел: частку і залишок.
	\begin{py}{Приклад}
x = 23
y = 7
print(divmod(x, y))   # (3, 2)
	\end{py}
	
	\subsection{Функція \texttt{hex(x)}}
	Конвертує ціле число у рядок із шістнадцяткових цифр.
	\begin{py}{Приклад}
x = 15
print(hex(x))   # '0xf'
	\end{py}
	
	\subsection{Функція \texttt{oct(x)}}
	Конвертує ціле число у рядок із вісімкових цифр.
	\begin{py}{Приклад}
x = 15
print(oct(x))   # '0o17'
	\end{py}
	
	\subsection{Функція \texttt{round(x)}}
	Округлює число до найближчого цілого.
	\begin{py}{Приклад}
print(round(3.0254))   # 3
	\end{py}
	
	\subsection{Функція \texttt{round(x, n)}}
	Округлює число \texttt{x} з точністю до \texttt{n} знаків після коми.
	\begin{py}{Приклад}
print(round(3.0254, 2))   # 3.03
	\end{py}
	